\subsection{16.38: periodic subgroups in a set product}
\label{cand:16.38}
\problemstatement{16.38}
\begin{proposition}
If \(G\) is soluble and \(A,B\le G\) are periodic, every subgroup
contained in the set \(AB\) is periodic.
\end{proposition}
\begin{proof}
First, periodic soluble groups are locally finite. A finitely generated
one has finite quotient by its last nontrivial derived subgroup,
by induction on derived length. That subgroup is finitely generated
by Schreier's theorem, and abelian periodic, hence finite.

We need a cocycle observation. Let a locally finite group \(C\) act on
a torsion-free abelian group \(D\), and let
\(\delta(cd)=\delta(c)+c\delta(d)\).
The image of \(\delta\) cannot contain both \(t\) and \(2t\) for \(t\ne0\).
Indeed, choose preimages \(c_1,c_2\), and put \(H=\langle c_1,c_2\rangle\),
which is finite. In \(\Q\otimes D\), averaging gives
\[
 v=|H|^{-1}\sum_{h\in H}\delta(h),\qquad
 \delta(h)=v-hv.
\]
Average an inner product on the finite-dimensional real span of \(Hv\).
The vectors \(v,v-t,v-2t\) then have equal norm. Expanding the last two
equalities yields
\(2\langle v,t\rangle=\|t\|^2\) and
\(4\langle v,t\rangle=4\|t\|^2\), so \(t=0\).
This averaging is over the finite group \(H\), not over all of \(C\).

Suppose now that \(D\trianglelefteq E\) is torsion-free abelian and
\(A,B\le E\) are locally finite. For \(\pi:E\to E/D\), let
\(C=\pi(A)\cap\pi(B)\). The preimages of \(C\) in \(A,B\) map
isomorphically to \(C\), since their intersections with \(D\) are trivial.
Write their inverse sections as \(a,b\). In
\(\pi^{-1}(C)=D\rtimes a(C)\), put \(b(c)=\delta(c)a(c)\);
then \(\delta\) is a cocycle and \(C\) is locally finite.
A factorization lying in \(D\) must have inverse quotient coordinates,
and hence
\[
 D\cap AB=\{a(c)b(c)^{-1}:c\in C\}=-\delta(C).
\]
The observation shows that this set cannot contain both \(g\) and
\(g^2\) for any \(1\ne g\in D\).

It suffices to prove that \(\langle g\rangle\subseteq AB\) implies
that \(g\) has finite order. Induct on the derived length of \(G\).
The abelian case is immediate. Otherwise let \(D\) be its last
nontrivial derived subgroup. Induction in \(G/D\) gives \(g^r\in D\)
for some \(r>0\). The torsion subgroup \(T\) of \(D\) is normal in \(G\),
and \(D/T\) is torsion-free abelian. Apply the preceding observation
in \(G/T\) to \(g^rT\) and \(g^{2r}T\), both of which belong to the
appropriate set product. It gives \(g^r\in T\), proving the claim.
\end{proof}

The set \(AB\) itself can contain elements of infinite order.
For example, let \(D=\bigoplus_{i\ge1}\Z\), let
\(C=\bigoplus_{i\ge1}C_2\) act by coordinate sign changes, and set
\(\delta(c)_i=c_i\) for \(c_i\in\{0,1\}\).
This is a cocycle. In \(D\rtimes C\), the two complements
\[
 A=\{(0,c)\},\qquad B=\{(\delta(c),c)\}
\]
are elementary abelian, whereas \(D\cap AB=-\delta(C)\) contains
nonzero infinite-order elements and excludes their doubles.
The theorem concerns subgroups contained in \(AB\).
It also does not assert the locally soluble version of Problem~9.65.
See \artifact{research/16.38-proof.md}.
